\section{The initial vector with common auxiliary labels}
\label{sec:replica_prevector}

Let $A$ be a finite set, $\Omega\in\mathbb C^A$, and $d,k\in\mathbb N$.
The initial vector below is the construction in
\cite[Section~7]{OpenAI2026AreaLaw}, equations
\textup{(comparator:prevector)} and \textup{(comparator:high-label)}.

\begin{definition}[Initial vector in product coordinates]
    \label{def:replica_prevector}
    \lean{TensorPower.replicaPrevector}
    \leanok
    Let $u_d=d^{-1/2}\sum_{i=0}^{d-1}|i,i\rangle$ be the uniform auxiliary pair
    for $d>0$, and set $u_0=0$. For $k=0$, use the empty tensor product.
    Regroup its literal tensor power into the two auxiliary systems and write it
    as $U_k$. For a central label projector $\Pi_\lambda$ on
    $(\mathbb C^d)^{\otimes k}$, define
    \begin{align}
        v_{k,\lambda}=\Omega^{\otimes k}\otimes
        (I_C\otimes\Pi_\lambda)U_k.
    \end{align}
    All products and projections here act on the actual finite copy spaces.
\end{definition}

\begin{theorem}[Symmetry of the explicit initial vector]
    \label{thm:replica_prevector_symmetry}
    \lean{TensorPower.replicaPrevector_copyPerm_fixed}
    \leanok
    \uses{def:replica_prevector}
    For every label $\lambda$ and every permutation $\sigma\in S_k$, let
    $R_A(\sigma)$ and $R_d(\sigma)$ permute the copies of the physical and
    auxiliary systems, respectively. Then
    \begin{align}
        (R_A(\sigma)\otimes(R_d(\sigma)\otimes R_d(\sigma)))v_{k,\lambda}
        =v_{k,\lambda}.
    \end{align}
\end{theorem}
\begin{proof}
    \leanok
    \uses{def:replica_prevector}
    Reindexing each repeated product gives
    \begin{align}
        R_A(\sigma)\Omega^{\otimes k}       & =\Omega^{\otimes k}, \\
        (R_d(\sigma)\otimes R_d(\sigma))U_k & =U_k.
    \end{align}
    The central projector $\Pi_\lambda$ commutes with $R_d(\sigma)$.
    Consequently $I_C\otimes\Pi_\lambda$ preserves the invariance of $U_k$,
    and tensoring with $\Omega^{\otimes k}$ gives the asserted identity.
\end{proof}

\begin{theorem}[The supplied label on both auxiliary systems]
    \label{thm:replica_prevector_auxiliary_labels}
    \lean{TensorPower.replicaPrevector_auxiliary_label_fixed}
    \leanok
    \uses{def:replica_prevector}
    Let $A$ be a finite set, let $\Omega\in\mathbb C^A$, and let
    $d,k\in\mathbb N$. For every label $\lambda$ of $S_k$, the initial
    vector $v_{k,\lambda}$ has label $\lambda$ on both auxiliary systems:
    \begin{align}
        (I_{A^k}\otimes(\Pi_\lambda\otimes I_R))v_{k,\lambda}
         & =v_{k,\lambda}, \\
        (I_{A^k}\otimes(I_C\otimes\Pi_\lambda))v_{k,\lambda}
         & =v_{k,\lambda}.
    \end{align}
    The identities include $d=0$, $k=0$, and the zero physical vector,
    with the conventions of Definition~\ref{def:replica_prevector}.
\end{theorem}
\begin{proof}
    \leanok
    \uses{def:replica_prevector,thm:uniform_bell_matching_labels}
    Put $w=(I_C\otimes\Pi_\lambda)U_k$. Idempotence gives
    $(I_C\otimes\Pi_\lambda)w=w$. By Theorem~\ref{thm:uniform_bell_matching_labels}, the central label
    projections have the same action on the two halves of the repeated uniform pair:
    $(\Pi_\lambda\otimes I_R)U_k=(I_C\otimes\Pi_\lambda)U_k$.
    Since the two projections commute, this also gives
    $(\Pi_\lambda\otimes I_R)w=w$. Tensoring both identities with
    $\Omega^{\otimes k}$ proves the assertion.
\end{proof}

Let $\Omega\in\mathbb C^A$ and
$\psi\in\mathbb C^S\otimes\mathbb C^d$ be unit vectors. Suppose that
$H\Omega=E_0\Omega$, with $E_0\in\mathbb R$.
Write $\rho=\operatorname{tr}_S|\psi\rangle\langle\psi|$.

\begin{theorem}[One sequence for the initial vector and selected mass]
    \label{thm:replica_prevector_sequence}
    \lean{TensorPower.exists_label_sequence_replicaPrevector}
    \leanok
    \uses{def:replica_prevector,thm:high_label_pure_sequence}
    There is one sequence $(\lambda_k)$ for which every $k\ge1$ satisfies
    \begin{align}
        0<\|v_{k,\lambda_k}\|\le1,
        \qquad (\overline H_k\otimes I_{CR})v_{k,\lambda_k}
        =E_0v_{k,\lambda_k},
        \qquad \overline H_k=\frac1k\sum_{i=0}^{k-1}H^{(i)}.
    \end{align}
    The vector has label $\lambda_k$ on both auxiliary systems and is invariant
    under every simultaneous copy permutation on the physical, left auxiliary
    and right auxiliary systems. For all sufficiently large $k$,
    \begin{align}
        \|(I_{S^k}\otimes\Pi_{\lambda_k})\psi^{\otimes k}\|^2
        \ge\frac1{2(k+1)^{d^2}},
    \end{align}
    and $\log d_{\lambda_k}=kS(\rho)+o(k)$.
    The initial projection of $U_k$ is required only to be nonzero; no inverse
    polynomial lower bound on its probability is asserted.
\end{theorem}
\begin{proof}
    \leanok
    \uses{def:replica_prevector,thm:high_label_pure_sequence}
    Start with the common sequence selected from the actual density $\rho$.
    For any remaining copy number at which the selected sector has zero mass,
    choose a sector of positive mass using the complete central resolution of
    $\rho^{\otimes k}$. The original sequence has positive mass eventually,
    so these changes affect only finitely many terms and preserve the mass
    bound and entropy asymptotic.

    Unit normalization of $\psi$ implies $d>0$. Every selected projector is
    nonzero, and hence acts nontrivially on the uniform pair. Tensoring with
    $\Omega^{\otimes k}$ preserves its norm. Orthogonal projection is contractive,
    and both unprojected tensor powers have norm one, giving the asserted norm
    bounds. The one-copy eigenvector equation gives the mean-energy equation
    on the physical product, irrespective of the auxiliary vector.

    The repeated uniform pair has matching left and right central labels.
    Idempotence therefore gives both fixed-label equations after projection.
    The uniform pair and the physical product are invariant under simultaneous
    copy permutations; centrality of the right label projector preserves this
    invariance after projection.
\end{proof}
